\section{Dyadic localization and its information cost}\label{l:dyadic}

We first prove the repeated-localization estimate that motivates the paper. The prior is the image of a uniform cube under a spherical coordinate map, not uniform surface measure on the patch.

Put
\[
 n=d-1,\qquad k=\lfloor n/16\rfloor,\qquad
 m=\lfloor n/8\rfloor,
\]
and assume $d$ is sufficiently large. Let $\mu_0$ be normalized Lebesgue
measure on the half-open cube
\[
 Q_0=[-1/(2\sqrt n),1/(2\sqrt n))^n,
 \qquad \phi(z)=(z,\sqrt{1-\|z\|^2}).
\]
Thus $\phi(Q_0)\subset S^{d-1}$. Half-open cells give each point a
unique cell at every level.

\begin{lemma}[Geometry of the cube chart]\label{ld:cube-geometry}
The map \(\phi\) is \(2\)-Lipschitz, and a level-\(J\) dyadic cell
has side \(2^{-J}/\sqrt n\) and \(\mu_0\)-mass \(2^{-nJ}\).
For every fixed \(v\in S^{d-1}\) and \(\epsilon>0\),
\[
 \mu_0\{z:\arccos\langle\phi(z),v\rangle\le\epsilon\}
 \le \bigl(\sqrt{2\pi e}\,\epsilon\bigr)^n
 \le(5\epsilon)^n.
\]
\end{lemma}
\begin{proof}
On \(Q_0\), \(\|z\|\le1/2\), so the gradient of
\(\sqrt{1-\|z\|^2}\) has norm at most \(1/\sqrt3\).
This proves the Lipschitz bound; the cell formulas follow from the
side length and normalized cube volume. Angular error at most
\(\epsilon\) implies chordal error at most \(\epsilon\), so the first
\(n\) coordinates lie in the radius-\(\epsilon\) ball about the
projection of \(v\). Its normalized cube volume is at most
\(n^{n/2}v_n\epsilon^n\). The Gaussian integral on the radius-\(\sqrt n\)
ball gives \(v_n\le(2\pi e/n)^{n/2}\), proving the claim.
\end{proof}

Draw $Z\sim\mu_0$. In block $i$, independently draw a standard Gaussian
$k\times d$ matrix $X_i$ and observe $(X_i,X_i\phi(Z))$. A message
$W_i$ takes at most $N$ values. Its conditional law may be any measurable
function of that exact block and the preceding messages; independent
randomization is allowed. More precisely, for each fixed preceding
message history the rule is an everywhere-defined Borel probability
kernel on $\R^{k\times d}\times\R^k$: its coordinates are nonnegative
and sum to one at every $(X,y)$. Euclidean spaces and the sphere carry
their Borel sigma fields. Neither the matrices nor their labels are
automatically retained in the message history.

Information and relative entropy use the conventions of Section~\ref{l:model}.

\begin{theorem}[Repeated localization]\label{ld:thm:main}
There is an absolute constant $c>0$ with the following property.
For any fixed $C_0>0$, consider the experiment above with $t\ge0$
blocks and $N\ge1$ possible message values, and suppose
\begin{equation}\label{ld:eq:alphabet}
 \bigl(2/m+e^{-ck}\bigr)\log N\le C_0d.
\end{equation}
One can reveal nested dyadic cells
$Q_0\supset Q_1\supset\cdots\supset Q_t$ containing $Z$, where $Q_i$
is a finite-valued measurable function of $Z$, the preceding revelations,
and $W_1,\ldots,W_i$, so that the augmented transcript
\[
 \Pi_t=(W_1,Q_1,\ldots,W_t,Q_t)
\]
satisfies
\begin{equation}\label{ld:eq:main}
 I(Z;\Pi_t)\le Cdt.
\end{equation}
Here $C$ is a constant depending only on $C_0$. If $J_t$ is the level of $Q_t$, the
construction also gives $\E J_t\le Ct$.
\end{theorem}

For fixed $A$, an alphabet of size at most $\exp(Ad^2)$ satisfies
\eqref{ld:eq:alphabet} with a constant depending only on $A$. The theorem
allows arbitrary rules inside a block; it restricts only what a block
passes to the next one. Since the original message history is a function
of $\Pi_t$, it inherits the same information bound.

The proof has three steps. Section~\ref{ld:sec:projection} shows that if a
posterior assigns little mass to every small dyadic cell, exact Gaussian
projection has controlled density even after an $L^2$ weight is inserted.
Section~\ref{ld:sec:message} turns that estimate into a bound on the
information in one finite message. Section~\ref{ld:sec:localization} selects
a deepest cell whose posterior mass is large relative to its scale and
bounds the entropy of that selection. The information chain rule then
closes the iteration in Section~\ref{ld:sec:iteration}.
Section~\ref{ld:sec:learners} gives a complete application to finite-state
learners, including bounded stopping.

\subsection{A weighted density for exact projections}\label{ld:sec:projection}

We first work inside one dyadic cell and put $\alpha=1/2$. The hypothesis below measures
concentration at every smaller scale. Its role is to prevent independent
points from accumulating near a low-dimensional affine subspace.

The weight in the coming estimate has a specific information-theoretic
purpose. A message's probability, after restricting labels to a bounded
ball, is a function of the signal. Testing that function against every
$L^2$ weight of norm at most one bounds its $L^2$ norm by duality.
Jensen's inequality then controls the message's relative-entropy
contribution using this likelihood norm. Section~\ref{ld:sec:message}
will carry out these two steps. We therefore need a projection-density
bound that is uniform over the weights, not just a density bound for
the unweighted posterior.

Let $Q$ be a level-$J$ cell and let $\mu$ be a probability measure
supported in $Q$. For $B\ge1$, assume
\begin{equation}\label{ld:eq:spread}
 \mu(R)\le B2^{-\alpha n l}
 \quad\text{for every level-$(J+l)$ descendant $R$ of $Q$ and $l\ge0$}.
\end{equation}
Write $r=2^{-J}$, $c_Q=\phi(\operatorname{center}Q)$, and
\[
 a(z)=\frac{\phi(z)-c_Q}{r}.
\]
Then $\|a(z)\|\le2$. Let $P_X$ denote the law of a standard Gaussian
$k\times d$ matrix.

\begin{lemma}[Weighted exact projection]\label{ld:lem:density}
Under \eqref{ld:eq:spread}, for every measurable $h$ with
$\|h\|_{L^2(\mu)}\le1$, put $\nu=|h|\mu$. The image of
$P_X(dX)\nu(dz)$ under $(X,z)\mapsto(X,Xa(z))$ has a density
$f_\nu$ relative to $P_X(dX)\,du$, where $du$ is Lebesgue measure
on $\R^k$, and
\begin{equation}\label{ld:eq:density}
 \|f_\nu\|_{L^m(P_Xdu)}\le C^d B^{1/2}
\end{equation}
for an absolute constant $C$.
\end{lemma}

\begin{proof}
We first derive a tube bound from \eqref{ld:eq:spread}, then estimate a
smoothed projection by expanding its $m$th moment. Finally we pass to
the exact measure. The replica expansion and Gaussian orthogonalization
have a close predecessor in \cite[Section~7, Lemma~12]{SSV2019}; the
weighted multirow estimate and the exact-measure passage are proved here.

For every affine $q$-plane $V\subset\R^d$ with $q\le m-2$, we claim
\begin{equation}\label{ld:eq:tube}
 \mu\{z:\dist(a(z),V)\le\delta\}
 \le BC^d\delta^{\alpha n-q}\qquad(0<\delta\le1).
\end{equation}
Projection onto the first $n$ coordinates shows that every such $z$
lies within $\delta r$ of the projection of $c_Q+rV$. Only the part of
that projected plane within $2r$ of the center of $Q$ is needed. A
maximal $\delta r$-separated subset of that part has at most
$(C/\delta)^q$ points, by volume comparison in the plane. Balls of
radius $2\delta r$ about these points cover all relevant $z$.

Choose $l\ge0$ so that $2^{-l}\le\delta<2^{-l+1}$. Each covering
ball meets at most $C^n$ level-$(J+l)$ cells. Indeed those cells are
disjoint, have side greater than $\delta r/(2\sqrt n)$, and lie
inside a ball of radius $3\delta r$. If $v_n$ is the volume of the
unit ball in $\R^n$, the estimate
$v_n\le(2\pi e/n)^{n/2}$ gives the claimed count; the powers of
$\sqrt n$ cancel. Each such cell has mass at most
$B2^{-\alpha nl}\le B\delta^{\alpha n}$, proving \eqref{ld:eq:tube}.

By Cauchy--Schwarz, $\nu$ has mass at most one. For $\eta>0$ define
\[
 F_\eta(X,u)=\int (2\eta)^{-k}
 \ind_{\{\|u-Xa(z)\|_\infty\le\eta\}}\,d\nu(z).
\]
This is smoothing in the label variable only. Expanding its $m$th
power introduces points $z_1,\ldots,z_m$. Put $a_i=a(z_i)$, define
$D_i$ for $2\le i\le m$, and set
\[
 D_i=\dist\bigl(a_i,\aff(a_1,\ldots,a_{i-1})\bigr),\qquad
 \mathsf D=[a_2-a_1,\ldots,a_m-a_1],\qquad
 V_*^2=\det(\mathsf D^{\mathsf T}\mathsf D).
\]
The exponent in \eqref{ld:eq:tube} is positive. Thus each affine span in
this display has zero $\mu$-mass and zero $\nu$-mass. Almost every
tuple is affinely independent, and Gram--Schmidt gives
$V_*=\prod_{i=2}^mD_i$.

The intersection of the $m$ smoothing boxes has volume at most
$(2\eta)^k$ and is empty unless
$\|X(a_i-a_1)\|_\infty\le2\eta$ for every $i\ge2$.
For a fixed affinely independent tuple, one row of these differences
is Gaussian with covariance $\mathsf D^{\mathsf T}\mathsf D$ and maximum density
$(2\pi)^{-(m-1)/2}V_*^{-1}$. The $k$ rows are independent.
Bounding the probability of the difference box by its volume times
this maximum density, and using Tonelli's theorem, yields
\begin{equation}\label{ld:eq:moment}
 \int F_\eta^m\,dP_X\,du
 \le C^{k(m-1)}\int V_*^{-k}\,d\nu^{\otimes m}.
\end{equation}
All powers of $\eta$ have canceled. We now bound the right-hand side
uniformly over the preceding points in each affine span.

Put $\gamma=\alpha n-(m-2)$. Cauchy--Schwarz and \eqref{ld:eq:tube}
give, for each fixed preceding tuple,
\[
 \nu\{D_i\le\delta\}\le(BC^d)^{1/2}\delta^{\gamma/2}
 \qquad(0<\delta\le1).
\]
The chosen $k,m$ satisfy $\gamma/2-k\ge n/8$. Integrating the tail
of $D_i^{-k}$ therefore gives
\begin{align*}
 \int D_i^{-k}\,d\nu
 &\le1+k\int_0^1\nu\{D_i\le\delta\}\delta^{-k-1}\,d\delta\\
 &\le1+\frac{k(BC^d)^{1/2}}{\gamma/2-k}
 \le C^dB^{1/2}.
\end{align*}
Integrate $z_m,z_{m-1},\ldots,z_2$ in this order in
\eqref{ld:eq:moment}; the remaining $z_1$ integral has mass at most one.
Taking the $m$th root gives
\[
 \|F_\eta\|_m
 \le\bigl[C^{k(m-1)}(C^dB^{1/2})^{m-1}\bigr]^{1/m}
 \le C^dB^{1/2}.
\]

To recover exact labels, let $\psi$ be continuous and compactly
supported in the variables $(X,u)$. Averaging $\psi$ over a shrinking
box and applying bounded convergence shows
\[
 \int\psi F_\eta\,dP_X\,du
 \longrightarrow\int\psi(X,Xa(z))\,dP_X\,d\nu(z).
\]
H\"older's inequality bounds this limiting functional by
$C^dB^{1/2}\|\psi\|_{m/(m-1)}$. Continuous compactly supported
functions are dense in $L^{m/(m-1)}(P_Xdu)$. Duality supplies an
$L^m$ density with that bound. Agreement on continuous compactly
supported tests identifies the two locally finite Borel measures.
This proves \eqref{ld:eq:density} for the exact joint law, so it can be
tested against arbitrary measurable message rules.
\end{proof}

\subsection{The information in one message}\label{ld:sec:message}

The weighted estimate allows us to control the probability of each
message as a function of the signal. That is the link between the
geometric hypothesis and information.

\begin{lemma}[One block]\label{ld:lem:block}
Let $\mu,Q,B$ satisfy \eqref{ld:eq:spread}. Draw $Z\sim\mu$ and a
standard Gaussian $k\times d$ matrix $X$ independently. Suppose $W$
has at most $N$ values and conditional probabilities
$\pi_w(X,X\phi(Z))$. Assume each $\pi_w$ is Borel on
$\R^{k\times d}\times\R^k$, with $\pi_w\ge0$ and
$\sum_w\pi_w(X,y)=1$ for every $(X,y)$.
Then
\begin{equation}\label{ld:eq:block}
 I_\mu(Z;W)\le Cd+\log B+\bigl(2/m+e^{-ck}\bigr)\log N
\end{equation}
for absolute constants $C,c>0$.
\end{lemma}

\begin{proof}
Keep the normalization $X\phi(z)=Xc_Q+rXa(z)$ from
Section~\ref{ld:sec:projection}. Let
$L=\{u\in\R^k:\|u\|\le4\sqrt k\}$. Since $\|a(z)\|\le2$,
\begin{equation}\label{ld:eq:tail}
 \Prob_X\{Xa(z)\notin L\}\le e^{-ck}
\end{equation}
uniformly in $z$. For example, if $G$ is standard Gaussian in
$\R^k$, then $\E e^{\|G\|^2/4}=2^{k/2}$, and Markov's inequality
at $\|G\|>2\sqrt k$ proves this with $c=1-\tfrac12\log2$.
Also $|L|=v_k(4\sqrt k)^k\le C^k$.

Define the subprobabilities
\[
 p_w(z)=\E_X\bigl[\pi_w(X,X\phi(z))\ind_{\{Xa(z)\in L\}}\bigr]
\]
For a reference experiment, keep $X$ Gaussian and replace the normalized
exact label $Xa(z)$ by an independent uniform point in $L$. Its message
probabilities are
\[
 b_w=|L|^{-1}\int_L\E_X\pi_w(X,Xc_Q+ru)\,du.
\]
They satisfy $\sum_wb_w=1$. This uses the probability-kernel identity
at every label, including labels outside the support of the observation
law. Test Lemma~\ref{ld:lem:density} with
$\pi_w(X,Xc_Q+ru)\ind_L(u)$. Because $0\le\pi_w\le1$,
H\"older's inequality gives, for every $\|h\|_2\le1$,
\[
 \int|h|p_w\,d\mu
 \le C^dB^{1/2}(|L|b_w)^{1-1/m}.
\]
Duality in $L^2(\mu)$ consequently gives
\begin{equation}\label{ld:eq:message-norm}
 \|p_w\|_2\le A b_w^{1-1/m},\qquad A=C^dB^{1/2}\ge1.
\end{equation}

We next turn this into information without assuming that the posterior
is uniform. Adjoin the bit $G=\ind_{\{Xa(Z)\in L\}}$ to $W$,
which can only increase mutual information. Put
$P_w=\int p_w\,d\mu$ and $P_g=\sum_wP_w$.
For $P_w>0$, Jensen's inequality under the probability measure
$p_w\mu/P_w$ gives
\[
 \int p_w\log(p_w/P_w)\,d\mu
 \le2P_w\log(\|p_w\|_2/P_w).
\]
Thus the contribution from $G=1$ to $I(Z;W,G)$ is at most
\begin{multline*}
 2\log A-(2-2/m)\sum_wP_w\log(P_w/b_w)\\
 {}+(2/m)\sum_wP_w\log(1/P_w).
\end{multline*}
If $b_w=0$, then \eqref{ld:eq:message-norm} gives $P_w=0$, and the
corresponding terms vanish. The log-sum inequality and the entropy
bound for $N$ outcomes give
\[
 \sum_wP_w\log(P_w/b_w)\ge P_g\log P_g,
 \qquad
 \sum_wP_w\log(1/P_w)\le\log N-P_g\log P_g.
\]
The good contribution is therefore at most
$2\log A+(2/m)\log N-2P_g\log P_g$.

For $G=0$, let $q_w(z)$ be the corresponding conditional
subprobability and $Q_w=\int q_w\,d\mu$. Since $q_w\le1$,
\[
 \sum_w\int q_w\log(q_w/Q_w)\,d\mu
 \le\sum_wQ_w\log(1/Q_w)
 \le P_b\log N-P_b\log P_b,
\]
where $P_b=1-P_g\le e^{-ck}$ by \eqref{ld:eq:tail}.
The functions $-P_g\log P_g$ and $-P_b\log P_b$ are bounded by
absolute constants. Finally $2\log A\le Cd+\log B$, proving
\eqref{ld:eq:block}.
\end{proof}

\subsection{Selecting a cell and paying for its entropy}\label{ld:sec:localization}

After a message, \eqref{ld:eq:spread} may fail with a useful value of $B$.
We restore it by selecting a cell using the current posterior. The
selection must be based on the event of choosing that cell, rather than
on the larger event that the signal merely lies in it.

Fix a finite-valued history $H$ and a value $h$ of probability $p>0$.
Suppose the posterior $\lambda=\mathcal L(Z\mid H=h)$ is supported
in a level-$J$ dyadic cell $Q$. A descendant $R$ at relative depth $j$
is called \emph{heavy} when
\[
 \lambda(R)\ge2^{-\alpha nj}.
\]
For $Z$ drawn from $\lambda$, let $U$ be its deepest heavy ancestor
inside $Q$, and write $j(U)$ for its relative depth.

\begin{lemma}[Deepest heavy cell]\label{ld:lem:heavy}
The selection $U$ is measurable and finite-valued. Put
$p_R=\Prob(U=R\mid H=h)$ for each selected cell of positive
probability, and define
\[
 B_R=\max\{1,2^{-\alpha n j(R)}/p_R\}.
\]
The conditional law $\mathcal L(Z\mid H=h,U=R)$ satisfies
\eqref{ld:eq:spread} with parent $R$ and constant $B_R$. Moreover,
\begin{align}
 H(U)&\le H(j(U))+\alpha n\log2\,\E j(U),\label{ld:eq:cell-entropy}\\
 H(j(U))&\le\log2\,(\E j(U)+1),\label{ld:eq:scale-entropy}\\
 \E\log B_U&\le H(j(U))+1.\label{ld:eq:excess-entropy}
\end{align}
Here every probability, entropy, and expectation is conditional on
$H=h$.
\end{lemma}

\begin{proof}
For every measurable set $A$, the posterior satisfies
$\lambda(A)\le\mu_0(A)/p$. A heavy cell at relative depth $j$
therefore satisfies
\[
 2^{-\alpha nj}\le p^{-1}2^{-n(J+j)}.
\]
Since $1-\alpha>0$, this bounds $j$. The parent cell is heavy, so
each path has a deepest heavy cell. There are only finitely many cells
at the relevant levels, and all selection events are measurable.

Fix a selected cell $R$ at relative depth $j$, and let $R'$ be a
strict descendant at further depth $l$. If $R'$ is heavy, the event
$\{U=R\}\cap R'$ is empty: any point of $R'$ has a deeper heavy
ancestor. If $R'$ is not heavy, then
$\lambda(R')<2^{-\alpha n(j+l)}$. In either case,
\[
 \Prob(Z\in R'\mid H=h,U=R)
 \le\frac{2^{-\alpha nj}}{p_R}2^{-\alpha nl}
 \le B_R2^{-\alpha nl}.
\]
At depth zero the bound follows from $B_R\ge1$. Notice that this
uses the probability $p_R$ of \emph{selection}; replacing $p_R$ by
$\lambda(R)$ would condition on a different event.

At depth $j$, there are at most $K_j=2^{\alpha nj}$ heavy cells,
because their posterior masses sum to at most one. The entropy chain
rule and the bound on the number of choices at each depth prove
\eqref{ld:eq:cell-entropy}. Comparing the depth distribution with the
geometric distribution $2^{-(j+1)}$, nonnegativity of relative entropy
gives \eqref{ld:eq:scale-entropy}.

For the remaining estimate, write $\log^+x=\max\{0,\log x\}$ for
$x>0$. Let $P_j=\Prob(j(U)=j)$ and, on a level
of positive mass, set $\rho_R=p_R/P_j$. There are at most $K_j$
possible cells on that level. Consequently, for $u\ge0$,
\[
 \Prob\!\left\{\log^+\frac1{K_j\rho_U}>u\,\middle|\,j(U)=j\right\}
 \le e^{-u}.
\]
Indeed every cell in this event has conditional mass less than
$e^{-u}/K_j$, and there are at most $K_j$ of them. Integrating the
tail shows that this positive logarithm has conditional expectation at
most one. Since
\[
 \log B_R
 =\log^+\frac1{K_jP_j\rho_R}
 \le\log(1/P_j)+\log^+\frac1{K_j\rho_R},
\]
averaging over the levels proves \eqref{ld:eq:excess-entropy}.
\end{proof}

The three estimates serve different purposes. The first pays for the
cell itself. The second pays for the scale used to name it. The third
bounds the extra concentration factor that enters the next application
of Lemma~\ref{ld:lem:block}. All three costs are governed by the same
expected increase in depth.

\subsection{Closing the information bound}\label{ld:sec:iteration}

We can now construct the cells in Theorem~\ref{ld:thm:main}. The proof
compares two bounds for information: an upper bound that charges the
messages and revelations, and a lower bound supplied by the volume of
the final cell.

\begin{proof}[Proof of Theorem~\ref{ld:thm:main}]
Start with $\Pi_0$ empty, $Q_0$ as above, $J_0=0$, and $B_0=1$.
After observing $W_i$, put $\Pi_i^-=(\Pi_{i-1},W_i)$. Apply
Lemma~\ref{ld:lem:heavy} to this history and the parent $Q_{i-1}$.
Let $Q_i$ be the selected cell, $j_i$ its relative depth, and
$J_i=J_{i-1}+j_i$. Set $\Pi_i=(\Pi_i^-,Q_i)$, and let $B_i$
be the constant supplied by that lemma.

Only positive-probability histories matter. Each history has a finite
range by induction: the message range is finite, and each intermediate
history permits finitely many selected cells by Lemma~\ref{ld:lem:heavy}.
Thus every entropy below is finite. Fresh rows remain independent of
$Z$ conditional on the past, because all past messages and revelations
depend only on $Z$, past data, and independent past randomization.
The rule for $W_i$, conditional on this enlarged history, is an admissible
kernel for Lemma~\ref{ld:lem:block}; the extra cells are analysis variables
and do not change the original rule.

Let $h_i=H(j_i\mid\Pi_i^-)$, including the average over intermediate
histories. The alphabet hypothesis and Lemmas~\ref{ld:lem:block}
and~\ref{ld:lem:heavy} give
\[
 I(Z;W_i\mid\Pi_{i-1})\le Cd+\E\log B_{i-1},
 \qquad \E\log B_i\le h_i+1.
\]
The selection is deterministic from $(Z,\Pi_i^-)$, so
\begin{align*}
 I(Z;Q_i\mid\Pi_i^-)
 &=H(Q_i\mid\Pi_i^-)\\
 &\le h_i+\alpha n\log2\,\E j_i,
 \qquad h_i\le\log2\,(\E j_i+1).
\end{align*}
Summing the information chain rule, using $B_0=1$ and
$\sum_i j_i=J_t$, yields
\begin{equation}\label{ld:eq:depth-upper}
 I(Z;\Pi_t)\le Cdt+(\alpha n+2)\log2\,\E J_t.
\end{equation}
The $2$ in the depth coefficient accounts for the depth entropy at
selection and for its use in the next block's concentration factor.

For a fixed terminal history, its posterior $\lambda$ is supported
in its level-$J_t$ cell $Q_t$. Writing $\mu_0(\,\cdot\mid Q_t)$
for normalized restriction to this cell gives the exact identity
\[
 \KL(\lambda\|\mu_0)
 =nJ_t\log2+
 \KL\bigl(\lambda\|\mu_0(\,\cdot\mid Q_t)\bigr).
\]
Nonnegativity and averaging imply
\begin{equation}\label{ld:eq:depth-lower}
 I(Z;\Pi_t)\ge n\log2\,\E J_t.
\end{equation}
Combining \eqref{ld:eq:depth-upper} and \eqref{ld:eq:depth-lower}, and
using $(1-\alpha)n-2\ge n/4$ for large $n$, proves
$\E J_t\le Ct$. Substitution into \eqref{ld:eq:depth-upper} proves
\eqref{ld:eq:main}.
\end{proof}

The strict inequality $\alpha<1$ is what makes this comparison close.
Revealing depth $J_t$ certifies $nJ_t\log2$ information, while the
heavy-cell count charges only $\alpha nJ_t\log2$, apart from the
two scale-entropy terms. No uniformity assumption on the posteriors is
used.

\subsection{Application to finite-state Gaussian regression}\label{ld:sec:learners}

The information theorem applies to the finite-state model of
Section~\ref{l:model}. This application uses the original every-signal
success premise to obtain success under the cube prior.

\begin{corollary}\label{ld:cor:streaming}
Suppose \(M(d)=o(d^2)\) and \(0<\epsilon(d)\le1/10\). If a family of
learners in the model of Section~\ref{l:model} satisfies
\[
 \Prob\{\arccos\langle\widehat s,s\rangle\le\epsilon(d)\}
 \ge2/3\qquad\text{for every }s\in S^{d-1},
\]
then \(T(d)\ge c d\log(1/\epsilon(d))\) for an absolute \(c>0\) and
all sufficiently large \(d\). The dimension threshold may depend on
the memory sequence.
\end{corollary}

\begin{proof}[Proof of Corollary~\ref{ld:cor:streaming}]
First average the success guarantee over $s=\phi(Z)$, $Z\sim\mu_0$.
Fix the shared rule-selection randomness, if present, with cube-prior
average success greater than $3/5$; such a choice exists because the
original average is at least $2/3$. The chart has full-rank derivative because its first \(n\) coordinates are \(z\), so Lemma~\ref{l:prior-ac} applies. Apply Lemma~\ref{l:borel} to these fixed rules. Now expose all remaining independent randomness and
fix a realization with prior-average success at least $3/5$.
Finite transition and stopping
kernels can be sampled using independent uniforms indexed by state and
sample index. Output laws can likewise be sampled in advance for each
terminal state and stopping index. The resulting deterministic learner
has the same state bound, and its output is a function of its terminal
state and stopping index.

Divide its horizon into $t=\lceil T/k\rceil$ blocks. Pregenerate all
$k$ rows in every block, including an unused final suffix. A block
message records one of the following: the continuing state; a terminal
state together with the local stopping offset; or a dummy value if the
learner was already inactive. Allowing offsets at the block endpoints
and immediate stopping, it suffices to use
\begin{equation}\label{ld:eq:stop-alphabet}
 N\le(k+3)2^M+1.
\end{equation}
The transcript determines the global stopping index from the block
index and local offset. Conditional on its past, the starting state is
known, and composing the transitions gives a measurable rule of the
whole exact block. Thus Theorem~\ref{ld:thm:main} applies. Eventually
$M\le d^2$, so \eqref{ld:eq:stop-alphabet} gives $\log N\le Cd^2$
and hence \eqref{ld:eq:alphabet} with an absolute $C_0$.

We next lower-bound the information needed for accurate estimation.
For a fixed unit vector $v$, angular error at most $\epsilon$ implies
$\|\phi(Z)-v\|\le\epsilon$, and hence
$\|Z-P_{\R^n}v\|\le\epsilon$. The normalized cube volume of this
event is at most
\begin{equation}\label{ld:eq:cap}
 n^{n/2}v_n\epsilon^n\le(5\epsilon)^n.
\end{equation}
Let $p\ge3/5$ be the success probability for the fixed learner and
the cube prior. Under the product of the marginal laws of
$(Z,\widehat s)$, the same success event has probability
$p_0\le(5\epsilon)^n$. Relative-entropy data processing to its binary
indicator gives
\begin{align*}
 I(Z;\widehat s)
 &\ge p\log(p/p_0)+(1-p)\log((1-p)/(1-p_0))\\
 &\ge p\log(1/p_0)-\log2\\
 &\ge pn\log(1/(5\epsilon))-\log2.
\end{align*}
For $\epsilon\le1/10$,
\[
 \log(1/(5\epsilon))
 \ge\left(1-\frac{\log5}{\log10}\right)\log(1/\epsilon).
\]
Since the output is determined by the messages, data processing and
Theorem~\ref{ld:thm:main} imply $t\ge c\log(1/\epsilon)$ for large $d$.

It remains to remove the ceiling in \(t=\lceil T/k\rceil\).
Return to the original randomized learner and average its every-signal
success over a uniform spherical signal. The \(T>d/8\) instance of
Lemma~\ref{l:caps-endpoint} applies to this same learner and prior; it
grants the entire stream and therefore includes early stopping. This
step does not infer cube-prior success from a sphere-average premise.
Thus \(T/k\ge1\), \(\lceil T/k\rceil\le2T/k\), and \(k\ge c'd\)
for large \(d\). The block lower bound gives
\(T\ge c''d\log(1/\epsilon)\), as required.
\end{proof}
